The construction of the Yogyakarta--Solo toll road may alter surface-road movement patterns in Yogyakarta, including along the Ring Road Utara (RRU) corridor. To support future evaluation of such network changes, a descriptive spatiotemporal baseline of observed vehicle movement is needed. This study analyzes trips indicated by devices classified as vehicles at six major intersections along RRU using active Mobile Positioning Data (MPD). The initial study-area dataset consists of 15,386,869 GPS points from 370,778 unique MAIDs. After R-tree-based spatial filtering, bilocation collapse, indicative motor-vehicle mode classification, map matching, and trajectory segmentation, the analysis dataset contains 212,407 pings from 8,382 MAIDs in 38,973 valid trips. The analysis focuses on intersection-zone origin-destination (OD) matrices, intersection intensity indicators, ping-per-trip distribution, and exploratory OD-zone movement patterns. The results indicate that 28,767 trips have identified intersection-zone OD pairs, with 58.48\% located on the main diagonal or O=D. The largest non-diagonal OD pair is Condongcatur--UPN with 956 trips. Intersection intensity indicators identify Condongcatur and Monjali as the zones with the highest average daily unique MAIDs, at 39.6 and 39.3 MAIDs/day, respectively. The trajectory distribution indicates sparse observations, with a median of 4 pings per trip and 70.37\% of trips containing at most 5 pings. Therefore, the results are interpreted as sample-based indicators of observed vehicles, not as actual traffic volume.

\textbf{Keywords:} mobile positioning data, Ring Road Utara, origin-destination, intersection intensity, sparse trajectory
